% ===== tr-02 块2：第 3 章 pp.11--14（原文页 11--14）=====
% [原文 p.11]
为了处理更有趣的情形，我们需要允许奇异纤维，从而允许 $B$ 的仿射结构中出现奇点。

\section{带奇点的仿射流形}

为了处理奇异纤维，我们作如下定义。

\begin{definition}{3.4.1}
（热带、整）带奇点的仿射流形是一个 $(C^0)$ 流形 $B$，它带有一个开子集 $B_0\subseteq B$，其上有（热带、整）仿射结构，并且使得 $\Gamma:=B\setminus B_0$ 是局部闭子流形的局部有限并，其余维数 $\ge2$。
\end{definition}

% [Tutor Note: 定义中 "(tropical, integral)" 与 "(C^0)" 的括号表示各修饰词可任取其一：即 B 只需是 C^0 流形，仿射结构可以是普通/tropical/integral 版本；奇点集 Γ 的余维数至少 2 是后文单值论证得以进行的关键（余维 1 会把 B_0 切断）。]

以例子说明：让我们解释 Batyrev 构造如何产生一大类这样的流形。该构造取自 [19]，其中对完全交给出了组合上更复杂的版本；另见 [28] 与 [29] 中的另一种构造。

设 $\Delta$ 是 $M_\R=M\otimes_\Z\R$ 中的自反多面体（reflexive polytope），其中 $M=\Z^n$；设 $N$ 为对偶格，$\nabla\subseteq N_\R$ 为如下给出的对偶多面体
\[
\nabla:=\{n\in N_\R \mid \langle m,n\rangle\ge-1\ \text{对所有 } m\in\Delta\}.
\]

我们假设 $0\in\Delta$ 是 $\Delta$ 的唯一内部格点。设 $\check\Sigma$ 为 $\nabla$ 的法扇，由 $\Delta$ 的各个面上的锥组成。假设给定 $\Delta$ 的一个星形细分，其所有顶点都是格点，并诱导出扇 $\check\Sigma$ 的一个细分 $\check\Sigma'$。此外假设
\[
\check h: M_\R\to\R
\]
是扇 $\check\Sigma'$ 上的（上）严格凸分段线性函数。又设
\[
\check\phi: M_\R\to\R
\]
为表示环面簇 $\PP_\nabla$ 的反典范类的分段线性函数；即 $\check\phi$ 在 $\check\Sigma$ 的每个一维锥的本原生成元上取值 1。最后，假设 $\check h$ 的选取使得 $\check h'=\check h-\check\phi$ 是一个（不一定严格的）凸函数。

对扇 $\check\Sigma'$ 上任意凸分段线性函数 $\check g$，定义 $\check g$ 的 Newton 多面体
\[
\nabla^{\check g}:=\{n\in N_\R \mid \langle m,n\rangle\ge-\check g(m)\ \text{对所有 } m\in M_\R\}.
\]
特别地，
\[
\nabla^{\check h}=\nabla^{\check h'}+\nabla^{\check\phi}=\nabla^{\check h'}+\nabla,
\]
其中 $+$ 表示 Minkowski 和。我们的目标是给 $B:=\partial\nabla^{\check h}$ 赋予一个带奇点的仿射结构。

我们实现这一目标的第一种方法不需要任何选择。

设 $\mathcal{P}$ 是 $\nabla^{\check h}$ 的各个真面构成的集。又设 $\operatorname{Bar}(\mathcal{P})$ 表示 $\mathcal{P}$ 的第一重心细分，并设 $\Gamma\subseteq B$ 为 $\operatorname{Bar}(\mathcal{P})$ 中所有不含 $\mathcal{P}$ 的顶点（一个零维胞腔）或者与 $\mathcal{P}$ 的某个极大胞腔的内部相交的单形之并。若令 $B_0:=B\setminus\Gamma$，我们可以如下在 $B_0$ 上定义仿射结构。$B_0$ 有一个开覆盖
\[
\{W_\sigma \mid \sigma\in\mathcal{P}\ \text{极大}\}\cup\{W_v \mid v\in\mathcal{P}\ \text{为顶点}\}
\]
% [Tutor Note: Γ 的取法是两类单形之并：「不含 P 的顶点（零维胞腔）的单形」与「与某个极大胞腔内部相交的单形」，两类都要；这样 Γ 的余维数才 ≥2，从而 B_0 是稠密开集、单值有定义。]
% [原文 p.12]
其中 $W_\sigma=\operatorname{Int}(\sigma)$ 是 $\sigma$ 的内部，而
\[
W_v=\bigcup_{\substack{\tau\in\operatorname{Bar}(\mathcal{P})\\ v\in\tau}}\operatorname{Int}(\tau)
\]
是 $v$ 在 $\operatorname{Bar}(\mathcal{P})$ 中的（开）星。我们定义仿射图
\[
\psi_\sigma: W_\sigma\to\mathbb{A}^{n-1}\subseteq N_\R
\]
由 $W_\sigma$ 到 $\mathbb{A}^{n-1}$（包含 $\sigma$ 的仿射超平面）的包含映射给出。再取
\[
\psi_v: W_v\to N_\R/\R v'
\]
为投影，其中 $v$ 作为 $\nabla^{\check h}$ 的顶点可以唯一地写成 $v'+v''$，这里 $v'$ 是 $\nabla$ 的顶点、$v''$ 是 $\nabla^{\check h'}$ 的顶点。容易验证：对 $v\in\sigma$，$\psi_\sigma\circ\psi_v^{-1}$ 是仿射线性的且其线性部分为整（整性来自 $\Delta$ 的自反性！），于是 $B$ 是带奇点的热带仿射流形。此外，若 $\check h$ 取为具有整斜率，则 $B$ 是整的。

% [Tutor Note: 记号 𝔸^{n−1} 是 (n−1) 维仿射空间（不是射影空间），它在这里等同于包含 σ 的仿射超平面；$\R^n\rtimes GL(\Z^n)$ 那类整线性部分来自 Δ 的自反性，这是 Batyrev 构造与镜像对称对偶性的组合基础。]

% [Tutor Note: ψ_v 的目标 $N_\R/\R v'$ 是模掉「由 v' 张成的实直线」所得的商空间；这里 v=v'+v'' 的分解唯一性依赖于 v 是 ∇^h 的顶点、v' 是 ∇ 的顶点、v'' 是 ∇^{h'} 的顶点。]

我们常常希望细化这一构造，以得到 $B$ 的更细的多面体分解 $\mathcal{P}$，并随之得到更有趣的判别轨迹 $\Gamma$。这样做的一个理由是：这一构造显然不是镜像对称的，因为它只依赖于 $\Delta$ 的星形细分而不依赖于 $\nabla$ 的。此外，$\nabla$ 的极大星形细分对应于 Batyrev 所称的 $\PP_\Delta$ 的 MPCP（极大投射 crepant 部分）消解；通常我们希望研究 $\PP_\Delta$ 的 MPCP 消解中的超曲面，而不是 $\PP_\Delta$ 本身中的超曲面。为引入这一额外的灵活性，我们需要作一些选择，具体如下。

首先，选取 $\nabla$ 的一个星形细分，使其所有顶点都是格点，并诱导出扇 $\Sigma$（即 $\Delta$ 的法扇）的加细 $\Sigma'$。这诱导出 $\partial\nabla$ 的一个多面体细分，我们把该细分的胞腔族记作 $\mathcal{P}_{\partial\nabla}$。注意，由于 $0\in\nabla$，有
\[
\nabla^{\check h'}\subseteq\nabla^{\check h'}+\nabla=\nabla^{\check h}.
\]

\begin{definition}{3.4.2}
若 $\partial\nabla^{\check h}$ 的细分 $\mathcal{P}$ 由 $\nabla^{\check h}$ 的一个细分 $\mathcal{P}_{\nabla^{\check h}}$ 诱导，且满足以下三条性质，则称 $\mathcal{P}$ 关于 $\mathcal{P}_{\partial\nabla}$ 是\emph{好的}：
\begin{enumerate}
\item $\nabla^{\check h'}$ 是 $\mathcal{P}_{\nabla^{\check h}}$ 中若干胞腔的并。
\item $\mathcal{P}_{\nabla^{\check h}}$ 的所有顶点都含于 $\partial\nabla^{\check h}$ 中或含于 $\nabla^{\check h'}$ 中。
\item 每个满足 $\sigma\cap\partial\nabla^{\check h}\ne\varnothing$ 且 $\tau:=\sigma\cap\partial\nabla^{\check h'}\ne\varnothing$ 的胞腔 $\sigma\in\mathcal{P}_{\nabla^{\check h}}$ 都可以写成
\[
\sigma=(C(\sigma')+\tau)\cap\nabla^{\check h},
\]
其中 $\sigma'\in\mathcal{P}_{\partial\nabla}$，而 $C(\sigma')$ 是 $\Sigma'$ 中对应的锥。
\end{enumerate}
若 $\check h$ 具有整斜率且 $\mathcal{P}_{\nabla^{\check h}}$ 的所有顶点都是整点，则称 $\mathcal{P}$ 是整的。
\end{definition}

% [原文 p.13]
下图展示了这样的好细分可能是什么样子，其中 $\nabla$ 是 $\mathcal{O}_{\PP^2}(3)$ 的 Newton 多面体：

\figimg[0.62]{p3_p13.png}

% [Tutor Note: 原文 p.13 中部此图无图号、无图注；译文按原文位置插入，不加 \figtitle、不自造图号（第 3 章的图均无编号）。]

% [Tutor Note: $\PP_\Delta$ 与 $\PP_\nabla$ 是两个对偶的环面簇（$\Delta$、$\nabla$ 互为对偶多面体的法扇所定义），MPCP 全称 maximal projective crepant partial resolution，即极大投射 crepant 部分消解；"partial" 不能漏译成「部分消解」之外的写法。]

给定 $B:=\partial\nabla^{\check h}$ 的一个好分解 $\mathcal{P}$，我们像之前一样在 $B$ 上再次得到带奇点的仿射结构：用这个新的多面体分解 $\mathcal{P}$ 的第一重心细分 $\operatorname{Bar}(\mathcal{P})$ 来定义判别轨迹 $\Gamma\subseteq B$。于是像之前一样，$B_0:=B\setminus\Gamma$ 有一个开覆盖
\[
\{W_\sigma \mid \sigma\in\mathcal{P}\ \text{极大}\}\cup\{W_v \mid v\in\mathcal{P}\ \text{为顶点}\}
\]
其中 $W_\sigma=\operatorname{Int}(\sigma)$ 是 $\sigma$ 的内部，而
\[
W_v=\bigcup_{\tau\in\operatorname{Bar}(\mathcal{P}),\,v\in\tau}\operatorname{Int}(\tau)
\]
是 $v$ 在 $\operatorname{Bar}(\mathcal{P})$ 中的（开）星。我们定义仿射图
\[
\psi_\sigma: W_\sigma\to\mathbb{A}^{n-1}\subseteq N_\R
\]
由 $W_\sigma$ 到 $\mathbb{A}^{n-1}$（包含 $\sigma$ 的仿射超平面）的包含映射给出。再取 $\psi_v: W_v\to N_\R/\R v'$ 为投影，其中 $v$ 可唯一地写成 $v'+v''$，$v'$ 是 $\nabla$ 的整点、$v''\in\nabla^{\check h'}$。像之前一样容易验证：对 $v\in\sigma$，$\psi_\sigma\circ\psi_v^{-1}$ 是仿射线性的且线性部分为整，于是 $B$ 是带奇点的热带仿射流形。此外，若 $\check h$ 有整斜率且 $\mathcal{P}$ 是整的，则 $B$ 上的仿射结构事实上是整的。

\begin{example}{3.4.3}
设 $\Delta\subseteq\R^4$ 是下列各点的凸包
\[
(-1,-1,-1,-1),\quad (1,0,0,0),\quad (0,1,0,0),\quad (0,0,1,0),\quad (0,0,0,1),
\]
于是 $\nabla$ 是下列各点的凸包
\[
(-1,-1,-1,-1),\quad (4,-1,-1,-1),\quad (-1,4,-1,-1),\quad (-1,-1,4,-1),\quad (-1,-1,-1,4).
\]
取 $\check h=\check\phi$ 并选取 $\nabla$ 的一个星形三角剖分。此时 $B=\partial\nabla$。容易看出 $B_0$ 上的仿射结构事实上可以越过 $\nabla$ 的所有三维面的内部延拓。这给出一个更小的判别轨迹 $\Gamma$；在 $\nabla$ 的一个好的正则三角剖分下，它在 $\nabla$ 的某个二维面的邻域中看起来像下图：细线给出三角剖分，粗线给出判别轨迹 $\Gamma$。
\end{example}

% [原文 p.14]
\figimg[0.62]{p3_p14.png}

% [Tutor Note: 原文 p.14 页首此图同样无图号、无图注；它承接 p.13 末句 "looks like the following picture in a neighbourhood of a 2-face of ∇"，译文按原文页序置于 p.14 起始处，不加图号。]

在这幅图中，$\Gamma$ 是一个三价图，带有两类三价顶点。沿着边缘的那些顶点是非平面的：图中画出的 $\Gamma$ 的两条附加腿含于 $\nabla$ 的另外两个二维面中。然而，在这个二维面的内部，$\Gamma$ 是平面的。

一般地，若 $\Sigma$ 与 $\check\Sigma$ 的细分 $\Sigma'$ 与 $\check\Sigma'$ 分别表示 $\PP_\Delta$ 与 $\PP_\nabla$ 的极大投射 crepant 部分消解，且 $\dim B=3$，那么只会出现这两类三价顶点。更具体地，人们知道局部系统 $\Lambda$ 与 $\check\Lambda$ 在 $B_0$ 上于这些顶点处的单值是什么样子。若 $v\in\Gamma$ 是含于某个二维面内部的顶点，则显然该二维面的切空间在 $\Lambda$ 绕 $v$ 的平行移动下不变。更细致的分析给出：在适当的基下，$\Lambda$ 的单值矩阵形如
\[
T_1=\begin{pmatrix}1&0&0\\1&1&0\\0&0&1\end{pmatrix},\quad
T_2=\begin{pmatrix}1&0&0\\0&1&0\\1&0&1\end{pmatrix},\quad
T_3=\begin{pmatrix}1&0&0\\-1&1&0\\-1&0&1\end{pmatrix}.
\]
这里 $T_1,T_2,T_3$ 由绕从 $v$ 伸出的 $\Gamma$ 的三条棱的回路作平行移动给出。当然，$\check\Lambda$ 的单值是这些矩阵的转置逆。类似地，若 $v$ 是含于 $\nabla^{\check h}$ 的某条棱中的 $\Gamma$ 的顶点，则单值将形如
\[
T_1=\begin{pmatrix}1&-1&0\\0&1&0\\0&0&1\end{pmatrix},\quad
T_2=\begin{pmatrix}1&0&-1\\0&1&0\\0&0&1\end{pmatrix},\quad
T_3=\begin{pmatrix}1&1&1\\0&1&0\\0&0&1\end{pmatrix}.
\]
于是我们看到，考察 $\Lambda$ 与 $\check\Lambda$，两类顶点的单值互相交换。

% [Tutor Note: 两组 T_i 都是 3×3 整矩阵；第一组对应「二维面内部」的顶点，第二组对应「nabla 上带 check 的 h 所定义的多面体 ∇ 的棱上」的顶点。文字层会把矩阵抽成线性字符串、甚至丢失负号，必须逐元素对照 p.14 页图（500 dpi 裁图复核过两组 9 个元素）。]

\begin{theorem}{3.4.4}
若 $B$ 是三维带奇点的热带仿射流形，使得 $\Gamma$ 是三价的，且 $\Lambda$ 在每个顶点处的单值是上述两类之一，则 $f_0:X(B_0)\to B_0$ 可以紧化为一个拓扑纤维化 $f:X(B)\to B$。对偶地，$\check f_0:\check X(B_0)\to B_0$ 可以紧化为一个拓扑纤维化 $\check f:\check X(B)\to B$。
\end{theorem}

% [Tutor Note: 定理 4.4 的结论是「拓扑纤维化」而非「特殊拉格朗日纤维化」：紧化只在拓扑（C^0）范畴内成立，这正是本节标题「带奇点的仿射流形」的拓扑进路的边界；不要译强。]
